% Homework: Decoding USB and CAN Frames.
% Do not compile this file directly; run Generate_Homework_USB_CAN.py, which
% generates the waveform figures and the answer data into ./build first.
% Compiling with \def\keymode{1} before \input produces the instructor key.
%
% By design, no PID table or field-format diagram is included here (unlike
% the in-class activity). Students label the stuff bits and field
% boundaries directly on each waveform figure, then answer a few short
% decode questions below it. They are expected to look up the relevant
% encoding rules, field formats, and PID codes in the USB 2.0 and classical
% CAN (ISO 11898-1) specifications -- that research is part of the
% assignment. CRC is not covered in this course, so no CRC questions are
% asked (the CRC field is still labeled as part of the frame structure).
\documentclass[11pt]{article}
\usepackage[landscape,margin=0.5in]{geometry}
\usepackage{graphicx}
\usepackage{enumitem}
\usepackage{xcolor}

\input{build/homework_data.tex}

\ifdefined\keymode
  \newcommand{\keyans}[1]{\par\smallskip{\color{red}\small\textbf{Key:} #1}\par}
  \newcommand{\docsuffix}{ -- Instructor Key}
\else
  \newcommand{\keyans}[1]{}
  \newcommand{\docsuffix}{}
\fi

% Ruled blank space students can use to write labels/notes if they would
% rather not mark on the figure itself.
\newcommand{\answerspace}[1]{\par\vspace{#1}\noindent\hrulefill\par\vspace{4pt}}

\setlength{\parindent}{0pt}
\setlist[enumerate]{leftmargin=*, itemsep=4pt, topsep=2pt}

\begin{document}

{\Large\bfseries ECGR 4101/5101 -- Advanced Embedded Systems}\hfill
Name: \underline{\hspace{5cm}}\quad Date: \underline{\hspace{2.5cm}}\\[2pt]
{\large Homework: Decoding USB and CAN Frames\docsuffix}

\smallskip
Each problem below gives you a waveform captured directly off the bus for one
frame. No reference tables or diagrams are provided -- you are expected to
look up the relevant encoding rules, field formats, and PID codes in the USB
2.0 specification (Problems 1-2) and the classical CAN / ISO 11898-1
specification (Problem 3). Finding this information yourself is part of the
assignment. \textbf{On each figure, label any stuff bit(s) you find and the
boundaries of each field (by bit time), then answer the short questions
below it.}

\newpage
\textbf{Problem 1: USB Full-Speed Packet}

\smallskip
The figure below is a capture of the D+ and D- lines during a single
\textbf{full-speed} USB data packet. The numbers along the top count the bit
times, and each dotted line marks the boundary between two bit times.

\begin{center}
  \includegraphics[width=\linewidth]{build/usb_full_waveform.png}
\end{center}

Label any stuff bit(s) and the boundaries of the SYNC, PID, DATA, CRC16, and
EOP fields directly on the figure above.
\keyans{Stuff bit(s) at bit time(s) \UsbFullStuffCells. SYNC
  \UsbFullSyncFirst--\UsbFullSyncLast, PID \UsbFullPidFirst--\UsbFullPidLast,
  DATA \UsbFullDataFirst--\UsbFullDataLast, CRC16
  \UsbFullCrcFirst--\UsbFullCrcLast, EOP \UsbFullEopFirst--\UsbFullEopLast.}

\answerspace{1.6cm}

\begin{enumerate}
  \item Decode the PID field. What type of packet is this?
    \underline{\hspace{4cm}}
    \keyans{\UsbFullPidName}
  \item What data byte(s) does this packet carry? Give your answer in hex.
    \underline{\hspace{5cm}}
    \keyans{\UsbFullDataBytes}
\end{enumerate}

\newpage
\textbf{Problem 2: USB Low-Speed Packet}

\smallskip
The figure below is a capture of the D+ and D- lines during a single
\textbf{low-speed} USB data packet. The numbers along the top count the bit
times, and each dotted line marks the boundary between two bit times.

\begin{center}
  \includegraphics[width=\linewidth]{build/usb_low_waveform.png}
\end{center}

Label any stuff bit(s) and the boundaries of the SYNC, PID, DATA, CRC16, and
EOP fields directly on the figure above.
\keyans{Stuff bit(s) at bit time(s) \UsbLowStuffCells. SYNC
  \UsbLowSyncFirst--\UsbLowSyncLast, PID \UsbLowPidFirst--\UsbLowPidLast,
  DATA \UsbLowDataFirst--\UsbLowDataLast, CRC16
  \UsbLowCrcFirst--\UsbLowCrcLast, EOP \UsbLowEopFirst--\UsbLowEopLast.}

\answerspace{1.6cm}

\begin{enumerate}
  \item Decode the PID field. What type of packet is this?
    \underline{\hspace{4cm}}
    \keyans{\UsbLowPidName}
  \item What data byte(s) does this packet carry? Give your answer in hex.
    \underline{\hspace{5cm}}
    \keyans{\UsbLowDataBytes}
\end{enumerate}

\newpage
\textbf{Problem 3: CAN Data Frame}

\smallskip
The figure below is a capture of the CAN\_H and CAN\_L lines during a single
\textbf{standard frame} (11-bit identifier) CAN data frame. The numbers
along the top count the bit times, and each dotted line marks the boundary
between two bit times. (Recall that for CAN, a dominant bit is a logical
\texttt{0} and a recessive bit is a logical \texttt{1}.)

\begin{center}
  \includegraphics[width=\linewidth]{build/can_waveform.png}
\end{center}

Label any stuff bit(s) and the boundaries of the SOF, identifier, RTR, IDE,
r0, DLC, DATA, CRC, CRC delimiter, ACK, and EOF fields directly on the
figure above. (Note that CAN's bit-stuffing rule is different from USB's.)
\keyans{Stuff bit(s) at bit time(s) \CanStuffCells. SOF
  \CanSofFirst--\CanSofLast, ID \CanIdFirst--\CanIdLast, RTR
  \CanRtrFirst--\CanRtrLast, IDE \CanIdeFirst--\CanIdeLast, r0
  \CanRsvdFirst--\CanRsvdLast, DLC \CanDlcFirst--\CanDlcLast, DATA
  \CanDataFirst--\CanDataLast, CRC \CanCrcFirst--\CanCrcLast, CRC delimiter
  \CanCrcDelimFirst--\CanCrcDelimLast, ACK \CanAckFirst--\CanAckLast, EOF
  \CanEofFirst--\CanEofLast.}

\answerspace{1.6cm}

\begin{enumerate}
  \item Decode the 11-bit identifier. Give your answer in hex.
    \underline{\hspace{4cm}}
    \keyans{\CanId}
  \item What data byte(s) does this frame carry? Give your answer in hex.
    \underline{\hspace{5cm}}
    \keyans{\CanDataBytes}
\end{enumerate}

\end{document}
